% GENERATED FILE. Edit canonical lessons, then run npm run generate:books.
% canonical-source-hash: fnv1a64:24db28570c8a0955
% canonical-lessons: GE-C386-klagen, GE-C386-klage, GE-C386-prozess, GE-C386-haften, GE-C386-haftung

\chapter{Vor Gericht}
\label{ch:ge-vor-gericht}
\hlchaptermodality{386}

\begin{chapteropening}
\textbf{By the end of this chapter:} \emph{I can talk about taking someone to court, a trial, and who is liable.}

Complete the last lesson of chapter 386: I can talk about taking someone to court, a trial, and who is liable.
\end{chapteropening}

\section[klagen]{klagen — to sue; to complain}

\label{lesson:GE-C386-klagen}



\begin{quote}
Before the new one: say the German for legally, then the German for the rule.
\end{quote}

\subsection*{You'll want to know: klagen}
\textbf{klagen} — \textquotedblleft{}to sue; to complain\textquotedblright{}.

\begin{quote}
Die Nachbarn klagen gegen den Lärm vom Flughafen.
\end{quote}

\begin{quote}
\emph{The neighbours are suing over the noise from the airport.}
\end{quote}

\emph{Gegen jemanden klagen} — to sue someone. In everyday speech: \emph{über Kopfschmerzen klagen} — to complain of a headache.

\begin{cousinweb}
\emph{klagen} is inherited from Proto-Germanic \emph{klagōną} (EtymDB), first \textquotedblleft{}to lament\textquotedblright{}; the legal sense, bringing your complaint before a court, grew out of it.
\end{cousinweb}

\begin{grammarlens}[title={What you've built}]
One more word for telling, explaining and describing at B1.
\end{grammarlens}

\subsection*{Your turn}
\begin{itemize}
\raggedright
  \item \emph{Say it:} \emph{klagen}
  \item \emph{Say it:} \emph{klagen}, once more
  \item \emph{Say it:} a sentence of your own with \emph{klagen}
  \item \emph{Recall:} say the German for legally, then the German for the rule, then say \emph{klagen} again
\end{itemize}

\begin{tcolorbox}[breakable,colback=teal!4,colframe=teal!35!black,title={Before you move on}]
What does klagen mean? (\textquotedblleft{}To sue; to complain\textquotedblright{}.) Say it once more.
\end{tcolorbox}
\section[die Klage]{die Klage — the lawsuit; the complaint}

\label{lesson:GE-C386-klage}



\begin{quote}
Before the new one: say the German for the rule, then the German for to sue.
\end{quote}

\subsection*{You'll want to know: die Klage}
\textbf{die Klage} — \textquotedblleft{}the lawsuit; the complaint\textquotedblright{}.

\begin{quote}
Das Gericht hat die Klage abgewiesen.
\end{quote}

\begin{quote}
\emph{The court dismissed the lawsuit.}
\end{quote}

Plural: \emph{die Klagen}. \emph{Klage erheben} — to file a lawsuit.

\begin{cousinweb}
\emph{Klage} is inherited from Middle High German \emph{klage} and Old High German \emph{klaga} \textquotedblleft{}a lament\textquotedblright{} (EtymDB), the noun of \emph{klagen}.
\end{cousinweb}

\begin{grammarlens}[title={What you've built}]
One more word for telling, explaining and describing at B1.
\end{grammarlens}

\subsection*{Your turn}
\begin{itemize}
\raggedright
  \item \emph{Say it:} \emph{die Klage}
  \item \emph{Say it:} \emph{die Klage}, once more
  \item \emph{Say it:} a sentence of your own with \emph{die Klage}
  \item \emph{Recall:} say the German for the rule, then the German for to sue, then say \emph{die Klage} again
\end{itemize}

\begin{tcolorbox}[breakable,colback=teal!4,colframe=teal!35!black,title={Before you move on}]
What does die Klage mean? (\textquotedblleft{}The lawsuit; the complaint\textquotedblright{}.) Say it once more.
\end{tcolorbox}
\section[der Prozess]{der Prozess — the trial; the process}

\label{lesson:GE-C386-prozess}



\begin{quote}
Before the new one: say the German for to sue, then the German for the lawsuit.
\end{quote}

\subsection*{You'll want to know: der Prozess}
\textbf{der Prozess} — \textquotedblleft{}the trial; the process\textquotedblright{}.

\begin{quote}
Der Prozess gegen den Politiker beginnt im Herbst.
\end{quote}

\begin{quote}
\emph{The trial of the politician starts in the autumn.}
\end{quote}

Plural: \emph{die Prozesse}. Also a process: \emph{ein langer Prozess}.

\begin{cousinweb}
\emph{Prozess} was borrowed into Middle High German as \emph{process} (EtymDB) from Medieval Latin \emph{prōcessus} \textquotedblleft{}a lawsuit\textquotedblright{}, Latin \textquotedblleft{}a going forward\textquotedblright{}, from \emph{prōcēdere} \textquotedblleft{}to go forward, to advance\textquotedblright{} (Lewis \& Short).
\end{cousinweb}

\begin{grammarlens}[title={What you've built}]
One more word for telling, explaining and describing at B1.
\end{grammarlens}

\subsection*{Your turn}
\begin{itemize}
\raggedright
  \item \emph{Say it:} \emph{der Prozess}
  \item \emph{Say it:} \emph{der Prozess}, once more
  \item \emph{Say it:} a sentence of your own with \emph{der Prozess}
  \item \emph{Recall:} say the German for to sue, then the German for the lawsuit, then say \emph{der Prozess} again
\end{itemize}

\begin{tcolorbox}[breakable,colback=teal!4,colframe=teal!35!black,title={Before you move on}]
What does der Prozess mean? (\textquotedblleft{}The trial; the process\textquotedblright{}.) Say it once more.
\end{tcolorbox}
\section[haften]{haften — to be liable; to stick}

\label{lesson:GE-C386-haften}



\begin{quote}
Before the new one: say the German for the lawsuit, then the German for the trial.
\end{quote}

\subsection*{You'll want to know: haften}
\textbf{haften} — \textquotedblleft{}to be liable; to stick\textquotedblright{}.

\begin{quote}
Eltern haften für ihre Kinder.
\end{quote}

\begin{quote}
\emph{Parents are liable for their children.}
\end{quote}

\emph{Für etwas haften} — to be liable for something. Also to stick: \emph{Der Aufkleber haftet nicht.}

\begin{cousinweb}
\emph{haften} goes back to Old High German \emph{haftēn} \textquotedblleft{}to cling, to be fastened\textquotedblright{}, usually traced to \emph{haft} \textquotedblleft{}held, captive\textquotedblright{}. EtymDB has the word but no history, so the account is hedged.
\end{cousinweb}

\begin{grammarlens}[title={What you've built}]
One more word for telling, explaining and describing at B1.
\end{grammarlens}

\subsection*{Your turn}
\begin{itemize}
\raggedright
  \item \emph{Say it:} \emph{haften}
  \item \emph{Say it:} \emph{haften}, once more
  \item \emph{Say it:} a sentence of your own with \emph{haften}
  \item \emph{Recall:} say the German for the lawsuit, then the German for the trial, then say \emph{haften} again
\end{itemize}

\begin{tcolorbox}[breakable,colback=teal!4,colframe=teal!35!black,title={Before you move on}]
What does haften mean? (\textquotedblleft{}To be liable; to stick\textquotedblright{}.) Say it once more.
\end{tcolorbox}
\section[die Haftung]{die Haftung — liability}

\label{lesson:GE-C386-haftung}



\begin{quote}
Before the new one: say the German for the trial, then the German for to be liable.
\end{quote}

\subsection*{You'll want to know: die Haftung}
\textbf{die Haftung} — \textquotedblleft{}liability\textquotedblright{}.

\begin{quote}
Für Jacken und Taschen übernehmen wir keine Haftung.
\end{quote}

\begin{quote}
\emph{We accept no liability for coats and bags.}
\end{quote}

\emph{GmbH} stands for \emph{Gesellschaft mit beschränkter Haftung} — limited liability company.

\begin{cousinweb}
\emph{Haftung} is the noun of \emph{haften} \textquotedblleft{}to be liable\textquotedblright{}, with \emph{-ung}; the verb goes back to Old High German \emph{haftēn} \textquotedblleft{}to cling\textquotedblright{}.
\end{cousinweb}

\begin{grammarlens}[title={What you've built}]
One more word for telling, explaining and describing at B1.
\end{grammarlens}

\subsection*{Your turn}
\begin{itemize}
\raggedright
  \item \emph{Say it:} \emph{die Haftung}
  \item \emph{Say it:} \emph{die Haftung}, once more
  \item \emph{Say it:} a sentence of your own with \emph{die Haftung}
  \item \emph{Recall:} say the German for the trial, then the German for to be liable, then say \emph{die Haftung} again
\end{itemize}

\begin{tcolorbox}[breakable,colback=teal!4,colframe=teal!35!black,title={Before you move on}]
What does die Haftung mean? (\textquotedblleft{}Liability\textquotedblright{}.) Say it once more.
\end{tcolorbox}
